\section{Pre-specified validation study}
\label{sec:protocol}

The analyses above use public data and simulation. A claim that BlinkTwin tracks real people needs a prospective study. We specify it here before collecting any data; it will be registered and will need approval from a human research ethics committee before recruitment.

\paragraph{Design.} Within-person, repeated-measures field study. Adults who work at a laptop for at least 4 hours a day use BlinkTwin on their own computer for 10 working days. On each day they enter their sleep, run the webcam during laptop work, and complete at least four check-ins, each preceded by a 3-minute reaction test, at prompted times spread across the day. A subset takes part in a supervised session of extended wakefulness to cover high sleepiness and to provoke microsleeps, subject to separate ethics approval and safety arrangements (for example, no driving home), with simultaneous EOG and video recording as ground truth.

\paragraph{Participants.} Healthy adults aged 18--60 without a diagnosed sleep disorder, without night-shift work during the study, and not taking medication that affects alertness. Informed consent covers on-device webcam processing and the export of de-identified event logs at the end of the study.

\paragraph{Hypotheses.} (H1) The anchored twin (variant E of Section~\ref{sec:simulation}) predicts held-out KSS ratings more accurately than the rhythm model alone (A) and than the population mappings without check-ins (C). (H2) It is not worse than check-ins alone (B), and its 90\% intervals are better calibrated than fusion without anchoring (D). (H3), exploratory: personal calibration helps in real people, that is, per-channel biases learned in the first days carry over to later days. Personal slope learning (F) is an exploratory comparison only, given its results in DROZY and in simulation.

\paragraph{Primary outcome.} Prediction error for held-out KSS ratings. For each check-in, the twin is replayed on all events before it, and its estimate is compared with the rating (leave-one-check-in-out, never using the future).

\paragraph{Secondary outcomes.} Prediction error for the next reaction-test median; Brier score and calibration of the forecast probability of KSS $\geq$ 7 over the next hour; error of the predicted time of reaching KSS~7 in the extended-wakefulness subset; sensitivity and false-alarm rate of the microsleep and head-drop alarms against EOG and video ground truth; agreement of webcam blink counts with EOG blink counts; usability (System Usability Scale \citep{brooke1996sus}); and adherence (completed check-ins, minutes of webcam coverage).

\paragraph{Analysis.} Absolute errors are analysed with linear mixed models with variant as a fixed effect and random intercepts for participant and day; leave-one-day-out cross-validation is used for any tuned parameter. Forecasts are evaluated with the Brier score, reliability diagrams and the area under the ROC curve.

\paragraph{Sample size.} The following are assumptions, not results. If the anchored twin reduces the mean absolute error of held-out KSS predictions by 0.3 points relative to the rhythm model alone (the simulation gave 0.305 with three check-ins), and the standard deviation of the per-person difference is 0.6 points, the within-person effect size is $d_z = 0.5$. A paired comparison at a two-sided $\alpha$ of 0.05 with 80\% power then needs about 34 participants. We plan to recruit 40 to allow for dropout. The much smaller difference between the anchored twin and check-ins alone in simulation ($-0.02$ to $-0.03$ KSS points) cannot be detected with this sample, which is why H2 is framed as non-inferiority and calibration rather than superiority. These assumptions will be updated from a pilot of 5 participants before the main study.
